\documentclass[a4paper, 12pt]{article}
\usepackage[T1]{fontenc}
\usepackage{amsmath}
\usepackage{algorithm}
\usepackage{algpseudocode}
\usepackage{textcomp}
\author{Karol Mamiński}
\title{Zestawienia Wybranych Komend Programu \LaTeX}
\date{}

\makeatletter
\newcommand{\linia}{\rule{\linewidth}{0.4mm}}
\renewcommand{\maketitle}{\begin{titlepage}
    \vspace*{1cm}
    \begin{center}\small
    Dokument wykonany\\
    w celach dydaktycznych
    \end{center}
    \vspace{2cm}
    \noindent\linia
    \begin{center}
      \LARGE \@title
         \end{center}
     \linia
    \vspace*{\stretch{6}}
    \begin{center}
    \@author
    \end{center}
  \end{titlepage}%
}
\makeatother

\renewcommand*\contentsname{Spis treści}

\renewcommand{\labelitemii}{\labelitemiv{}}

\begin{document}
\maketitle
\tableofcontents
\newpage
\section{Podstawy Składni}
\subsection{Komendy}
Komendy w \LaTeX zaczynają się znakiem $\backslash$ (backslash), po którym następują słowa kluczowe, np.:\newline
$\backslash$newline\newline
Wiele komend posiada również atrybuty, zarówno wymagane jak i opcjonalne. Pierwsze z wymienionych umieszczamy w nawiasach klamrowych (\{\}), drugie zaś w kwadratowych ([ ]). Przykład:\newline
$\backslash$documentclass[ a4paper, 12pt ]\{article\}

\subsection{Środowiska}
Poza standardowymi komendami do naszej dyspozycji w \LaTeX 'u posiadamy także środowiska. W prostych słowach są to komendy przeznaczone do stosowania na wielu wierszach.\newline
Środowiska rozpoczynamy za pomocą komendy $\backslash$begin\{\} i kończymy za pomocą $\backslash$end\{\}. Przykład:\newline
$\backslash$begin\{center\}\newline
Zawartość podlegająca formatowaniu środowiska. W tym przypadku wyśrodkowaniu. \\
$\backslash$end\{center\}

\subsection{Komendy własne}
Zdarza się, że materiały dostarczone przez \LaTeX 'a nie odpowiadają w pełni naszym potrzebom. Być może nie odpowiada nam domyślne formatowanie danej komendy, albo męczy nas wpisywanie za każdym razem dodatkowych atrybutów? W obu tych przypadkach możemy skorzystać z możliwości zmiany komendy bądź zdefiniowania własnej.\\
\\
W tym celu należy zastosować polecenia $\backslash$newcommand oraz $\backslash$renewcommand. Pierwsze z wymienionych pozwala nam na zdefiniowanie nowego polecenia, zaś druga na przypisanie nowych działań do już istniejącej komendy.\\
\\
Rozważmy dwa przypadki. W pierwszym chcemy mieć możliwość szybkiego nałożenia na tekst kursywy, podkreślenia i pogrubienia bez konieczności wpisywania trzech oddzielnych komend za każdym razem. Jednocześnie nie chcemy zmieniać działania poszczególnych komend, ponieważ są nam one potrzebne indywidualnie w innych częściach dokumentu. \\
W tym przypadku należy zastosować polecenie $\backslash$newcommand w następujący sposób:\\
$\backslash$newcommand\{wyroznienie\}[ 1 ]\{$\backslash$textbf\{$\backslash$emph\{$\backslash$underline\{\#1\}\}\}\}\\
$\backslash$wyroznienie\{tekst który chcemy wyróżnić\}\newline 
[ 1 ] -- ustalamy ile parametrów przyjmuje nasza komenda\\
\#1 -- odwołanie do pierwszego podanego parametru\\
\\
W~drugim przypadku chcemy zmienić sposób numerowania wszystkich list w dokumencie na małe litery rzymskie.\\
W~takim wypadku należy użyć polecenia $\backslash$renewcommand w~następujący sposób:\\
$\backslash$renewcommand\{$\backslash$labelenumi\}\{$\backslash$roman\{enumi\}\} \\


\section{Podział Dokumentu}
\subsection{Preambuła}
W preambule umieszczane są informacje kluczowe z punktu widzenia programu, takie jak klasa dokumentu, wybrane pakiety oraz dane, które mają znaleźć się w tytule.\newline
Najczęściej wykorzystywane na zajęciach elementy preambuły przedstawiono poniżej.
\begin{itemize}
	\item $\backslash$documentclass[ a4paper, 12pt ]\{article\} -- ustawia klasę dokumentu na artykuł, rozmiar papieru na A4 oraz rozmiar czcionki na 12 punktów
	\item $\backslash$usepackage[ T1 ]\{fontenc\} -- pozwala na pracę w języku polskim
	\item $\backslash$usepackage\{amsmath\} -- zapewnia dodatkowe opcje matematyczne
	\item $\backslash$usepackage\{algorithm\} -- pozwala na zapisywanie algorytmów
	\item $\backslash$usepackage\{algpseudocode\} -- pozwala na zapisywanie algorytmów
	\item $\backslash$author\{Imię i Nazwisko\} -- ustawia dane autora
	\item $\backslash$title\{Tytuł\} -- ustawia tytuł
	\item $\backslash$date\{data\} -- ustawia datę do wyświetlenia pod tytułem
\end{itemize}

\subsection{Treść dokumentu}
Właściwą treść dokumentu należy umieścić wewnątrz otoczenia \textbf{document}. \newline
Przykład:\newline
$\backslash$begin\{document\}\newline
Treść dokumentu\newline
$\backslash$end\{document\}
\newpage

\section{Formatowanie Tekstu}
\subsection{Wcięcia i odstępy}
\LaTeX rozpoznaje spacje oraz złamania wierszy w kodzie źródłowym jako pojedynczy odstęp. \\
Do dokładniejszej kontroli wcięć i odstępów wykorzystujemy następujące polecenia:
\begin{itemize}
	\item $\backslash$newline lub \begin{math}\backslash\backslash\end{math} -- łamie wiersz w wybranym miejscu (przechodzi do nowej linii)
	\item \begin{math}\backslash\backslash\end{math}* -- zabrania podziału strony w miejscu złamania wiersza
	\item $\backslash$newpage -- wstawia podział strony 
	\item $\backslash$clearpage -- rozpoczyna nową stronę
	\item $\backslash$cleardoublepage -- rozpoczyna nową nieparzystą stronę (w razie potrzeby tworzona również pusta parzysta strona)
	\item $\sim$ (tylda) -- wstawia odstęp zabraniając jednocześnie złamania wiersza (np. W\~{}poniedziałek)
\end{itemize} 

\subsection{Podział hierarchiczny}
W \LaTeX 'u dysponujemy następującymi znakami podziału:
\begin{itemize}
	\item w klasie article:
	\begin{itemize}
		\item $\backslash$section\{\}
		\item $\backslash$subsection\{\}
		\item $\backslash$subsubsection\{\}
		\item $\backslash$paragraph\{\}
		\item $\backslash$subparagraph\{\}
		\item $\backslash$appendix\{\}
	\end{itemize}
	\item w klasach book i report dodatkowo:
	\begin{itemize}
		\item $\backslash$chapter\{\}
		\item $\backslash$part\{\} -- nie narusza numeracji
	\end{itemize}
\end{itemize}
W przypadku każdej z w/w komend umieszczenie gwiazdki (*) tuż przed nawiasem klamrowym spowoduje brak numeracji (np. $\backslash$section*\{\} ).

\subsection{Znaki specjalne}
Znaki specjalne dostępne w \LaTeX :
\begin{itemize}
	\item $\backslash$\~{}\{\} -- \~{}
	\item \$$\backslash$sim\$ -- $\sim$
	\item $\backslash$ldots -- \ldots
	\item \$30$^{}$\{$\backslash$circ\}$\backslash$mathrm\{C\}\$ -- $30^{\circ}\mathrm{C}$
	\item $\backslash$texteuro -- \texteuro ~(wymaga pakietu textcomp)
	\item $\backslash$LaTeX -- \LaTeX
	\item aby uzyskać długi myślnik (taki jak powyżej) należy wpisać dwa minusy
\end{itemize}
\end{document}













